
\section{Introduction}

\subsection{Motivation}

Transformer-based language models concentrate the majority of their parameters 
in two architectural blocks: the feed-forward network and the multi-head 
self-attention (MHSA) mechanism. In DistilBERT-base, the four attention 
projections $W_Q, W_K, W_V, W_O$ account for approximately \textbf{70\%} of the 
parameters in each transformer block. Reducing the parameter cost of attention 
is therefore a high-leverage compression strategy.

Standard MHSA stores four dense matrices per layer:

\[
W_Q, W_K, W_V, W_O \in \mathbb{R}^{d \times d}
\tag{1.1}
\]

totaling $4d^2$ parameters. For DistilBERT ($d = 768$), this is $4 \times 589{,}824 = 2{,}359{,}296$ 
parameters per layer -- or 14.2M across six layers, a substantial share of the 
66.9M total.

\textbf{Low-rank factorization} is the classical approach. For each $W$, one 
replaces it with a product:

\[
W \approx U V, \qquad
U \in \mathbb{R}^{d_\text{out} \times r}, \quad
V \in \mathbb{R}^{r \times d_\text{in}}, \quad r \ll \min(d_\text{in}, d_\text{out})
\tag{1.2}
\]

The parameter count drops from $d_\text{in} d_\text{out}$ to 
$r(d_\text{in} + d_\text{out})$. The \textbf{reduction factor} is 
$\rho = d_\text{in} d_\text{out} / (r(d_\text{in} + d_\text{out}))$. For square 
matrices at $d = 768$ and $r = 128$: $\rho = 768 / 256 = 3.0$. At $r = 64$: 
$\rho = 6.0$. Classical SVD (Eckart--Young, 1936) gives the optimal rank-$r$ 
approximation in Frobenius norm.

\textbf{The rank-selection problem.} Classical SVD requires choosing $r$ 
a-priori, which is a hard trade-off:

\begin{itemize}
\item \textbf{Too small} -- the approximation error is high, and downstream accuracy
  drops.
\item \textbf{Too large} -- the parameter reduction is minimal, defeating the purpose.

\end{itemize}
In practice, different layers have different \textbf{intrinsic ranks}. Attention 
projections in early layers of DistilBERT tend to have rapidly-decaying 
spectra (effective rank ~30--48), while middle layers decay more slowly 
(effective rank ~80--100). A single global $r$ is suboptimal for both.

\subsection{Problem Statement}

Given a pretrained transformer $\mathcal{M}$ with $N_l$ layers and a target 
energy threshold $\eta \in (0, 1)$, we seek a compressed model 
$\mathcal{M}'$ such that:

\[
\mathcal{M}' \approx \mathcal{M}, \qquad
\frac{|\theta(\mathcal{M}')|}{|\theta(\mathcal{M})|} \ll 1
\tag{1.3}
\]

subject to:

\begin{enumerate}
\item \textbf{Per-layer adaptive rank.} Different layers may use different
  $r_\ell$ chosen automatically.
\item \textbf{Spectral energy preservation.} For each factorized layer,
  $E(r_\ell) \geq \eta$.
\item \textbf{Parameter reduction.} $r_\ell (d_\text{in} + d_\text{out}) < d_\text{in} d_\text{out}$.
\item \textbf{Multi-architecture support.} Flat (DistilBERT), nested (RoBERTa),
  causal with Conv1D (GPT-2), GQA + RoPE (Qwen).
\item \textbf{No base retraining.} Only LoRA adapters may be trained.

\end{enumerate}
\subsection{Our Approach: MinimalPRISM}

\textbf{MinimalPRISM} is a drop-in replacement for the four attention projections 
with three key components:

\begin{enumerate}
\item \textbf{SVD factorization.} Each $W$ is factorized as $U V$ via truncated SVD,
  with $U \in \mathbb{R}^{d_\text{out} \times r_\ell}$ and
  $V \in \mathbb{R}^{r_\ell \times d_\text{in}}$, where $\Sigma^{1/2}$ is
  distributed symmetrically.

\item \textbf{Spectral energy-adaptive rank.} For each weight matrix, compute
  $E(r) = \sum_{i \leq r} \sigma_i^2 / \sum_i \sigma_i^2$ and select:

  $$
  r\_\ell = \min\!\left(r\_\text{max}, \left|\{i : E(i) < \eta\}\right| + 1\right)
  \tag{1.4}
  $$

  This selects the smallest rank that captures at least $\eta$ fraction of
  the spectral energy, capped at $r_\text{max}$.

\item \textbf{LoRA recovery.} After factorization, the base factors $U, V$ are
  frozen and LoRA adapters $A, B$ are added:

  $$
  y = (x V^\top) U^\top + (x A^\top) B^\top \cdot s + b
  \tag{1.5}
  $$

  with $B = 0$ at initialization, so $y$ at step 0 equals the SVD
  approximation.

\end{enumerate}
\textbf{Injection safety.} MinimalPRISM detects nested encoders (RoBERTa) at 
injection time rather than only in weight-copy. It runs a **pre-test on the 
first layer** before mutating the rest of the model, aborting if weight 
extraction fails. This is critical for production correctness.

\subsection{Contributions}

\begin{itemize}
\item \textbf{Energy-adaptive rank selection.} A principled criterion that selects
  a different rank per layer based on spectral decay, avoiding the
  fixed-rank trade-off.
\item \textbf{A safe injection protocol.} Nested-encoder detection at injection time,
  pre-test on the first layer, and strict Frobenius norm validation of
  extracted weights.
\item \textbf{A multi-architecture implementation.} Four families supported
  (DistilBERT, RoBERTa, GPT-2, Qwen3.5) with explicit Conv1D handling, GQA
  support, and RoPE integration.
\item \textbf{A systematic evaluation.} Benchmark on three models on SST-2 with
  ablations over $r_\text{max}$, $\eta$, and LoRA rank.
\item \textbf{Full reproducibility.} Code, checkpoints, and scripts released under
  Apache 2.0.

\end{itemize}
\subsection{Paper Organization}

Section 2 reviews related work. Section 3 provides background. Section 4 
presents MinimalPRISM. Section 5 describes architecture integration. Section 
6 details the setup. Section 7 reports results. Section 8 analyzes 
energy-adaptive behavior. Section 9 discusses guidance. Section 10 lists 
limitations. Section 11 concludes.



\subsection{Notation}

\begin{lstlisting}

---
\end{lstlisting}